\section{线性表}

\subsection{线性表的定义与ADT}

\begin{mydef}{线性表}{linear-list}
线性表（Linear List）是由 $n\ (n\geqslant 0)$ 个数据元素 $a_1,a_2,\ldots,a_n$ 构成的有限序列，满足：
\begin{enumerate}
    \item 除首元素 $a_1$ 外，每个元素有且仅有一个直接前驱；
    \item 除末元素 $a_n$ 外，每个元素有且仅有一个直接后继；
    \item 元素个数 $n$ 称为表长，$n=0$ 时为空表。
\end{enumerate}
线性表是逻辑结构概念，不规定存储方式。实现时可选择顺序存储或链式存储。
\end{mydef}

\begin{conclusion}{线性表的抽象数据类型（ADT）}{adt-linear-list}
线性表 ADT 的核心操作集（以 C 风格描述）：
\begin{tablebox}{线性表 ADT 操作}
\centering
\begin{tabular}{|c|l|}
\hline
操作 & 语义 \\
\hline
\texttt{InitList(\&L)} & 初始化空表 \\
\texttt{DestroyList(\&L)} & 销毁表，释放空间 \\
\texttt{ClearList(\&L)} & 清空为初始状态 \\
\texttt{ListEmpty(L)} & 判空，返回 bool \\
\texttt{ListLength(L)} & 返回元素个数 \\
\texttt{GetElem(L, i, \&e)} & 取第 $i$ 个元素（$1\leqslant i\leqslant n$）\\
\texttt{LocateElem(L, e)} & 定位第一个等于 $e$ 的元素位置 \\
\texttt{PriorElem(L, cur, \&pre)} & 取前驱 \\
\texttt{NextElem(L, cur, \&next)} & 取后继 \\
\texttt{ListInsert(\&L, i, e)} & 在第 $i$ 位前插入 $e$ \\
\texttt{ListDelete(\&L, i, \&e)} & 删除第 $i$ 位元素，用 $e$ 返回 \\
\texttt{ListTraverse(L)} & 遍历访问 \\
\hline
\end{tabular}
\end{tablebox}

408 考试中重点关注 \texttt{GetElem}、\texttt{LocateElem}、\texttt{ListInsert}、\texttt{ListDelete} 的算法实现与复杂度，且要能区分类似 \texttt{LocateElem} 与 \texttt{GetElem} 的差异（按值 vs 按位）。
\end{conclusion}

\subsection{顺序表（Sequential List）}

\begin{mydef}{顺序表}{seq-list}
用一组\textbf{地址连续}的存储单元依次存储线性表的数据元素。若每个元素占 $s$ 个存储单元，首地址为 $\mathtt{LOC}(a_1)$，则第 $i$ 个元素的地址为：
\[
\mathtt{LOC}(a_i) = \mathtt{LOC}(a_1) + (i-1)\times s,\qquad 1\leqslant i \leqslant n.
\]
这一性质称为\textbf{随机存取}（Random Access）：任意位置元素的访问时间为 $O(1)$。
\end{mydef}

\begin{formula}{顺序表的存储结构（C 实现）}{seq-list-struct}
\begin{lstlisting}[language=C]
#define MaxSize 100          // 最大容量
typedef int ElemType;        // 元素类型，可按需替换
typedef struct {
    ElemType data[MaxSize];  // 存储空间
    int length;              // 当前长度
} SqList;                    // 顺序表类型
\end{lstlisting}
\textbf{关键点：} \texttt{data} 是静态数组，预先分配 \texttt{MaxSize}；\texttt{length} 记录实际元素个数（不是最后一个元素的下标）。408 常考「\texttt{length} 与下标的关系」——第 $i$ 个元素对应 \texttt{data[i-1]}。
\end{formula}

\subsubsection{顺序表的基本操作}

\begin{conclusion}{顺序表的插入操作}{seq-list-insert}
在第 $i$（$1\leqslant i \leqslant n+1$）位插入新元素 $e$：
\begin{enumerate}
    \item 检查 $i$ 是否合法，以及表是否已满；
    \item 将第 $i$ 到第 $n$ 个元素依次后移一位；
    \item 在位置 $i-1$ 处放入 $e$，\texttt{length++}。
\end{enumerate}

\begin{lstlisting}[language=C]
bool ListInsert(SqList *L, int i, ElemType e) {
    if (i < 1 || i > L->length + 1) return false;
    if (L->length >= MaxSize)        return false;
    for (int j = L->length; j >= i; j--)
        L->data[j] = L->data[j-1];
    L->data[i-1] = e;
    L->length++;
    return true;
}
\end{lstlisting}

\textbf{时间复杂度：}
\[
T(n) = \frac{1}{n+1}\sum_{i=1}^{n+1}(n-i+1) = \frac{n}{2}.
\]
等概率下平均移动 $n/2$ 个元素，为 $O(n)$。最好情况（插在末尾）$O(1)$，最坏情况（插在表头）$O(n)$。
\end{conclusion}

\begin{conclusion}{顺序表的删除操作}{seq-list-delete}
删除第 $i$（$1\leqslant i \leqslant n$）位元素并用 $e$ 返回：

\begin{lstlisting}[language=C]
bool ListDelete(SqList *L, int i, ElemType *e) {
    if (i < 1 || i > L->length) return false;
    *e = L->data[i-1];
    for (int j = i; j < L->length; j++)
        L->data[j-1] = L->data[j];
    L->length--;
    return true;
}
\end{lstlisting}

等概率下平均移动 $(n-1)/2$ 个元素，为 $O(n)$。
\end{conclusion}

\begin{conclusion}{顺序表按值查找}{seq-list-locate}
\begin{lstlisting}[language=C]
int LocateElem(SqList L, ElemType e) {
    for (int i = 0; i < L.length; i++)
        if (L.data[i] == e) return i + 1;  // 返回位序
    return 0;  // 未找到
}
\end{lstlisting}
平均比较次数 $(n+1)/2$，时间复杂度 $O(n)$。

\textbf{与按位查找的区别：} \texttt{GetElem(L, i, \&e)} 直接通过下标访问 \texttt{L.data[i-1]}，时间复杂度 $O(1)$——顺序表的随机存取特性。
\end{conclusion}

\subsection{单链表（Singly Linked List）}

\begin{mydef}{单链表}{linked-list}
用一组\textbf{任意}的存储单元存放线性表元素，每个结点包含：
\begin{itemize}
    \item \textbf{数据域}（data）：存放元素值；
    \item \textbf{指针域}（next）：指向直接后继结点的指针。
\end{itemize}
$n$ 个结点通过指针链成一个链表。最后一个结点的指针域为 \texttt{NULL}。

单链表不再支持随机存取——访问第 $i$ 个元素需从头沿链走 $i-1$ 步，时间复杂度 $O(n)$。
\end{mydef}

\begin{formula}{单链表的结点定义}{linked-list-node}
\begin{lstlisting}[language=C]
typedef struct LNode {
    ElemType data;
    struct LNode *next;
} LNode, *LinkList;
\end{lstlisting}
\texttt{LinkList} 是指向 \texttt{LNode} 的指针，即头指针。后续代码中的 \texttt{L} 通常指头指针。
\end{formula}

\subsubsection{头结点 vs 无头结点}

\begin{conclusion}{头结点的作用}{head-node}
408 考试中单链表通常\textbf{带头结点}（head node），即第一个结点（头结点）的 \texttt{data} 无实际意义，\texttt{next} 指向首元结点。好处：
\begin{enumerate}
    \item 空表/非空表的操作代码统一——头指针 \texttt{L} 始终非空（至少指向头结点）；
    \item 链表头部的插入/删除不需要单独判定边界。
\end{enumerate}
\begin{center}
\begin{tikzpicture}[>=stealth, node distance=1.8cm, auto]
\node[draw, rectangle] (H) {头结点}; 
\node[draw, rectangle, right of=H] (A) {$a_1$};
\node[draw, rectangle, right of=A] (B) {$a_2$};
\node[right of=B] (dots) {$\cdots$};
\node[draw, rectangle, right of=dots] (Z) {$a_n$};
\node[draw, rectangle, right of=Z] (N) {NULL};
\draw[->] (H) -- (A);
\draw[->] (A) -- (B);
\draw[->] (B) -- (dots);
\draw[->] (dots) -- (Z);
\draw[->] (Z) -- (N);
\node[above of=H, yshift=-0.8cm] {L（头指针）};
\draw[->] (0.3,0.8) -- (H);
\end{tikzpicture}
\end{center}
\textbf{注意：} 题目既可能采用带头结点的链表，也可能采用不带头结点的链表，务必以题设为准。
\end{conclusion}

\subsubsection{单链表的基本操作}

\begin{conclusion}{头插法建表}{head-insert}
每次将新结点插入到\textbf{头结点之后}，形成逆序链表：
\begin{lstlisting}[language=C]
LinkList CreateListHead(int n) {
    LinkList L = (LinkList)malloc(sizeof(LNode));
    L->next = NULL;                     // 头结点
    for (int i = 0; i < n; i++) {
        LNode *s = (LNode*)malloc(sizeof(LNode));
        scanf("%d", &s->data);
        s->next = L->next;              // （1） s 后继 = 原首元
        L->next = s;                    // （2） 头$\to$s
    }
    return L;
}
\end{lstlisting}
时间复杂度 $O(n)$。输入顺序与结果顺序相反。
\end{conclusion}

\begin{conclusion}{尾插法建表}{tail-insert}
每次将新结点接到链表\textbf{末尾}，保持插入顺序：
\begin{lstlisting}[language=C]
LinkList CreateListTail(int n) {
    LinkList L = (LinkList)malloc(sizeof(LNode));
    LNode *r = L;                       // r 始终指向当前尾结点
    for (int i = 0; i < n; i++) {
        LNode *s = (LNode*)malloc(sizeof(LNode));
        scanf("%d", &s->data);
        r->next = s;                    // 尾结点指向新结点
        r = s;                          // 更新尾指针
    }
    r->next = NULL;
    return L;
}
\end{lstlisting}
时间复杂度 $O(n)$，输入与结果顺序相同。
\end{conclusion}

\begin{conclusion}{单链表按位查找}{linked-list-get-elem}
从第一个\textbf{数据结点}（不是头结点）开始计数：
\begin{lstlisting}[language=C]
LNode *GetElem(LinkList L, int i) {
    if (i < 0) return NULL;
    LNode *p = L;                // p 指向头结点（第0个）
    int j = 0;
    while (p && j < i) {         // j=i 时 p 指向第i个结点
        p = p->next;
        j++;
    }
    return p;                    // 若 i>表长 则返回 NULL
}
\end{lstlisting}
\textbf{注意：}\texttt{GetElem(L, 0)} 返回头结点；\texttt{GetElem(L, 1)} 返回首元结点。时间复杂度 $O(n)$。
\end{conclusion}

\begin{conclusion}{单链表插入操作}{linked-list-insert}
在第 $i-1$ 个结点后插入新结点（即插入为第 $i$ 个）：
\begin{lstlisting}[language=C]
bool ListInsert(LinkList L, int i, ElemType e) {
    LNode *p = GetElem(L, i-1);       // 找第 i-1 个结点
    if (!p) return false;
    LNode *s = (LNode*)malloc(sizeof(LNode));
    s->data = e;
    s->next = p->next;                // （1） 新结点先接后继
    p->next = s;                      // （2） 前驱再接新结点
    return true;
}
\end{lstlisting}
\textbf{关键：} 语句（1）（2）的顺序不能颠倒！若先执行 \texttt{p->next = s}，则原后继链丢失。时间复杂度：查找 $O(n)$，插入本身 $O(1)$。
\end{conclusion}

\begin{conclusion}{单链表删除操作}{linked-list-delete}
删除第 $i$ 个结点：
\begin{lstlisting}[language=C]
bool ListDelete(LinkList L, int i, ElemType *e) {
    LNode *p = GetElem(L, i-1);       // 找第 i-1 个结点（前驱）
    if (!p || !p->next) return false; // p不存在或p无后继
    LNode *q = p->next;               // q 指向第 i 个结点
    *e = q->data;
    p->next = q->next;                // 前驱跳过 q
    free(q);
    return true;
}
\end{lstlisting}
\end{conclusion}

\begin{attention}
\textbf{408 常考陷阱——删除指定结点 p 本身：} 若只给指向 p 的指针（不知道前驱），可以将 p 的后继结点的 data 复制到 p，然后删除后继结点。但此方法\textbf{不适用于 p 是最后一个结点}，也不适用于结点 data 量过大或包含不能简单复制的资源（如指针/文件句柄）。
\end{attention}

\subsection{双链表（Doubly Linked List）}

\begin{mydef}{双链表}{double-linked-list}
每个结点有两个指针域：\texttt{prior} 指向前驱，\texttt{next} 指向后继。
\begin{lstlisting}[language=C]
typedef struct DNode {
    ElemType data;
    struct DNode *prior, *next;
} DNode, *DLinkList;
\end{lstlisting}
双链表可以从任意结点出发沿两个方向遍历，但存储密度低于单链表（多一个指针域）。
\end{mydef}

\begin{conclusion}{双链表的插入操作}{dlink-insert}
在结点 \texttt{p} 之后插入新结点 \texttt{s}：
\begin{lstlisting}[language=C]
s->next = p->next;                   // （1） s接p的后继
s->prior = p;                        // （2） s接p
if (p->next) p->next->prior = s;     // （3） p后继的前驱指向s
p->next = s;                         // （4） p接s
\end{lstlisting}
四条语句缺一不可，且顺序有讲究：先建立 s 的两个指针，再修改原链中指向 s 的两个指针。若 p 为尾结点则跳过（3）。
\end{conclusion}

\begin{conclusion}{双链表的删除操作}{dlink-delete}
删除结点 \texttt{p} 的后继结点 \texttt{q}：
\begin{lstlisting}[language=C]
DNode *q = p->next;
p->next = q->next;                   // p跳过q
if (q->next) q->next->prior = p;     // q后继的前驱指向p
free(q);
\end{lstlisting}
注意 \texttt{q} 不存在和 \texttt{q} 为尾结点的边界处理。
\end{conclusion}

\subsection{循环链表（Circular Linked List）}

\begin{mydef}{循环链表}{circular-list}
将单链表（或双链表）的最后一个结点的指针域指向头结点（而不是 NULL），形成一个环。
\begin{itemize}
    \item \textbf{循环单链表}：尾结点 \texttt{next} 指向头结点。判空条件：\texttt{L->next == L}。
    \item \textbf{循环双链表}：头结点的 \texttt{prior} 指向尾结点；尾结点的 \texttt{next} 指向头结点。判空条件：\texttt{L->next == L \&\& L->prior == L}。
\end{itemize}
循环链表适合需要\textbf{循环遍历}或从任意结点出发沿链反复扫描的场景。它避免了到达表尾后重新回到表头的特殊处理，但从一个任意结点查找到另一个任意结点通常仍需 $O(n)$ 时间，并不提供随机访问。
\end{mydef}

\subsection{静态链表（Static Linked List）}

\begin{mydef}{静态链表}{static-list}
用\textbf{数组}模拟链表：数组元素即结点，包含 \texttt{data} 和 \texttt{next}（记录后继结点的数组下标，-1 表示 NULL）。
\begin{lstlisting}[language=C]
#define MaxSize 100
typedef struct {
    ElemType data;
    int next;            // 后继下标
} SLinkList[MaxSize];
\end{lstlisting}
静态链表不需要动态内存分配（\texttt{malloc/free}），适用于不支持指针的语言或内存受限场景。408 中偶尔以选择题形式考察。
\end{mydef}

\subsection{顺序表 vs 链表的比较}

\begin{tablebox}{顺序表与链表的对比}
\centering
\begin{tabular}{|l|c|c|}
\hline
比较维度 & 顺序表 & 链表 \\
\hline
存储方式 & 连续地址 & 任意地址，指针链接 \\
存取方式 & 随机存取 $O(1)$ & 顺序存取 $O(n)$ \\
插入/删除 & 需移动元素 $O(n)$ & 修改指针 $O(1)$（已定位时）\\
存储密度 & $\approx 1$（纯数据） & $<1$（含指针域）\\
缓存性能 & 好（局部性好） & 差（结点分散）\\
动态扩容 & 静态数组需整体搬移 & 可按需分配单个结点 \\
\hline
\end{tabular}
\end{tablebox}

\subsection{408 高频考点与真题风格}

\begin{conclusion}{线性表 408 选择题常考点}{ds-chap1-keypoints}
\begin{enumerate}
    \item \textbf{顺序表插入/删除的元素移动次数}：平均 $\approx n/2$，注意区分插入与删除的差别。
    \item \textbf{头结点与无头结点}：空表判定、边界处理的代码差异。
    \item \textbf{单链表插入语句的顺序}：先连后继，再断前驱链。
    \item \textbf{「删除 p 所指结点」陷阱}：无前驱时可否用复制后继法、何时失效。
    \item \textbf{双链表/循环链表判空条件}：\texttt{L->next == L} 等。
    \item \textbf{时间复杂度辨析}：已定位结点的插入/删除是 $O(1)$，但\textbf{定位本身}需要时间——单链表按位查找 $O(n)$，顺序表按值查找 $O(n)$。
    \item \textbf{静态链表}：\texttt{next} 域存的是下标而非指针，增删元素不移动数据元素。
\end{enumerate}
\end{conclusion}

\newpage
\section{习题}

\subsection{基础过关题}

\begin{enumerate}

\item \textbf{（填空题）} 在一个长度为 $n$ 的顺序表中删除第 $i$（$1 \leqslant i \leqslant n$）个元素，等概率下需要向前移动的元素个数的期望值为 \underline{\qquad\qquad}；若在第 $i$（$1 \leqslant i \leqslant n+1$）个位置之前插入一个新元素，等概率下需要向后移动的元素个数的期望值为 \underline{\qquad\qquad}。

\ifshowsolutions\par\noindent\textbf{解析：} 删除第 $i$ 个元素需移动 $n-i$ 个元素，等概率下期望为 $\frac{1}{n}\sum_{i=1}^{n}(n-i)=\frac{n-1}{2}$。在第 $i$ 个位置前插入需移动 $n-i+1$ 个元素，等概率（$n+1$ 种可能）下期望为 $\frac{1}{n+1}\sum_{i=1}^{n+1}(n-i+1)=\frac{n}{2}$。

\boxed{\text{答案：}\dfrac{n-1}{2}\text{；}\dfrac{n}{2}}\fi

\item \textbf{（单项选择题）} 在单链表中，要将指针 \texttt{s} 所指的新结点插入到指针 \texttt{p} 所指结点之后，正确的语句序列是
\begin{center}
\begin{tabular}{ll}
(A) \texttt{s->next = p->next; p->next = s;} \\[4pt]
(B) \texttt{p->next = s; s->next = p->next;} \\[4pt]
(C) \texttt{s->next = p; p->next = s;} \\[4pt]
(D) \texttt{p->next = s->next; s->next = p;}
\end{tabular}
\end{center}

\ifshowsolutions\par\noindent\textbf{解析：} 插入结点时必须先让新结点 \texttt{s} 指向原后继 \texttt{p->next}，再修改 \texttt{p} 的指针域指向 \texttt{s}。若顺序颠倒（先执行 \texttt{p->next = s}），则原后继链丢失，无法再找到。故 (A) 正确。(B) 先断链导致丢失后继；(C) 是插在 \texttt{p} 之前而非之后；(D) 逻辑错误。\boxed{\textbf{答案：}(A)}\fi

\item \textbf{（真题风格·综合题）} 设单链表的结点结构为 \texttt{[data | next]}，头指针为 \texttt{L}（带头结点）。编写算法，将链表就地逆置，即辅助空间复杂度为 $O(1)$。分析算法的时间复杂度与空间复杂度。

\ifshowsolutions\par\noindent\textbf{解析：}
采用头插法思想：将头结点与后续结点断开，然后依次将每个数据结点以头插法重新插入到头结点之后。

\begin{lstlisting}[language=C]
void Reverse(LinkList L) {
    LNode *p = L->next;    // p 指向首元结点
    LNode *q;
    L->next = NULL;        // 头结点与后续断开，作为新链表的头
    while (p != NULL) {
        q = p->next;       // 暂存后继
        p->next = L->next; // 头插：p 的后继 = 原首元
        L->next = p;       // 头结点指向新首元 p
        p = q;             // p 后移
    }
}
\end{lstlisting}

\textbf{时间复杂度：} $O(n)$，遍历链表一次。\\
\textbf{空间复杂度：} $O(1)$，仅使用常数个辅助指针。

\textbf{关键点：} （1） 先断开 \texttt{L->next = NULL}，避免形成环；（2） \texttt{q} 暂存后继是必需的，否则修改 \texttt{p->next} 后无法找到下一个结点；（3） 头插法保证逆序结果。
\fi

\item \textbf{（真题风格·综合题）} 已知一个带有头结点的单链表，结点结构为 \texttt{[data | next]}。假设该链表只给出了头指针 \texttt{L}。在不改变链表的前提下，设计一个尽可能高效的算法，查找链表中倒数第 $k$（$k$ 为正整数）个位置上的结点。若查找成功，算法输出该结点的 \texttt{data} 值并返回 1；否则只返回 0。

\ifshowsolutions\par\noindent\textbf{解析：}
采用双指针（快慢指针）法，一趟扫描完成：

\begin{lstlisting}[language=C]
int SearchKthFromEnd(LinkList L, int k) {
    LNode *fast = L->next, *slow = L->next;
    int count = 0;
    // fast 先走 k 步
    while (fast != NULL && count < k) {
        fast = fast->next;
        count++;
    }
    if (count < k) return 0;   // k 超过表长
    // fast 和 slow 同步前进，fast 到 NULL 时 slow 即倒数第 k 个
    while (fast != NULL) {
        fast = fast->next;
        slow = slow->next;
    }
    printf("%d\n", slow->data);
    return 1;
}
\end{lstlisting}

\textbf{时间复杂度：} $O(n)$，只遍历一次链表。\\
\textbf{空间复杂度：} $O(1)$。

\textbf{思想：} 维护两个指针，间距为 $k$。当快指针走到表尾时，慢指针恰好指向倒数第 $k$ 个结点。
\fi

\item \textbf{（真题风格·综合题）} 已知一个带头结点的单链表，结点结构为 \texttt{[data | next]}。设计一个算法，删除链表中所有数据值在区间 $[a,b]$ 内的结点（$a \leqslant b$），并分析算法的时间复杂度。

\ifshowsolutions\par\noindent\textbf{解析：}
一趟遍历，维护前驱指针。由于单链表删除结点需要前驱，用 \texttt{p->next} 指向当前考察结点。

\begin{lstlisting}[language=C]
void DeleteRange(LinkList L, int a, int b) {
    LNode *p = L;            // p 指向待删结点的前驱
    LNode *q;
    while (p->next != NULL) {
        if (p->next->data >= a && p->next->data <= b) {
            q = p->next;
            p->next = q->next;   // 前驱跳过被删结点
            free(q);
        } else {
            p = p->next;          // 不删则后移
        }
    }
}
\end{lstlisting}

\textbf{时间复杂度：} $O(n)$，遍历链表一次。\\
\textbf{空间复杂度：} $O(1)$。

\textbf{关键：} （1） 用 \texttt{p->next} 考察结点，避免维护前驱的复杂性；（2） 删除时 \texttt{p} 不动（因为新的 \texttt{p->next} 可能也需删除），不删除时 \texttt{p} 才后移；（3） 注意区间是否包含端点，必须严格按题设处理开、闭区间。
\fi

\item \textbf{（真题风格·综合题）} 假定采用带头结点的单链表保存单词。当两个单词有相同的后缀时，可共享相同的后缀存储空间。设 \texttt{str1} 和 \texttt{str2} 分别指向两个单词所在单链表的头结点，链表结点结构为 \texttt{[data | next]}。设计一个尽可能高效的算法，找出由 \texttt{str1} 和 \texttt{str2} 所指的两个链表共同后缀的起始位置。

\ifshowsolutions\par\noindent\textbf{解析：}
先计算两链表表长，将较长链表的指针先移动 $|len1-len2|$ 步，使两指针到表尾的距离相等；然后同步后移，第一个相同结点即为公共后缀起点。

\begin{lstlisting}[language=C]
int ListLen(LinkList L) {
    int len = 0;
    LNode *p = L->next;
    while (p != NULL) { len++; p = p->next; }
    return len;
}

LNode *FindCommonSuffix(LinkList str1, LinkList str2) {
    int len1 = ListLen(str1), len2 = ListLen(str2);
    LNode *p = str1->next, *q = str2->next;
    // 将较长链表先走 |len1-len2| 步
    if (len1 > len2)
        for (int i = 0; i < len1 - len2; i++) p = p->next;
    else
        for (int i = 0; i < len2 - len1; i++) q = q->next;
    // 同步后移，找第一个相同结点
    while (p != NULL && p != q) {
        p = p->next;
        q = q->next;
    }
    return p;  // NULL 表示无公共后缀
}
\end{lstlisting}

\textbf{时间复杂度：} $O(\max(m,n))$，$m$、$n$ 分别为两表表长。\\
\textbf{空间复杂度：} $O(1)$。
\fi

\end{enumerate}

\subsection{真题风格综合训练}

\begin{enumerate}[resume]

\item \textbf{（真题风格·综合题）} 设将 $n$（$n>1$）个整数存放到一维数组 $R$ 中。设计一个在时间和空间两方面都尽可能高效的算法，将 $R$ 中保存的序列循环左移 $p$（$0<p<n$）个位置，即将 $R$ 中的数据由 $(x_0,x_1,\ldots,x_{n-1})$ 变换为 $(x_p,x_{p+1},\ldots,x_{n-1},x_0,x_1,\ldots,x_{p-1})$。

\ifshowsolutions\par\noindent\textbf{解析：}
核心思想：三次逆置。
\begin{enumerate}
    \item 逆置前 $p$ 个元素：$(x_{p-1},\ldots,x_1,x_0,\ldots)$；
    \item 逆置后 $n-p$ 个元素：$(\ldots,x_{n-1},\ldots,x_p)$；
    \item 整体逆置：$(x_p,\ldots,x_{n-1},x_0,\ldots,x_{p-1})$。
\end{enumerate}

\begin{lstlisting}[language=C]
void Reverse(int R[], int left, int right) {
    while (left < right) {
        int temp = R[left];
        R[left] = R[right];
        R[right] = temp;
        left++; right--;
    }
}

void LeftShift(int R[], int n, int p) {
    if (p <= 0 || p >= n) return;
    Reverse(R, 0, p - 1);       // 逆置前 p 个
    Reverse(R, p, n - 1);       // 逆置后 n-p 个
    Reverse(R, 0, n - 1);       // 整体逆置
}
\end{lstlisting}

\textbf{时间复杂度：} $O(n)$，每个元素参与两次交换（整体逆置时一次，分段逆置时一次），共约 $n$ 次交换。\\
\textbf{空间复杂度：} $O(1)$，仅使用常数个辅助变量。

\textbf{关键：} 三次逆置法是本题最优解，比逐位移动（$O(np)$）和辅助数组（$O(n)$ 空间）都更优。
\fi

\item \textbf{（真题风格·综合题）} 一个长度为 $L$（$L \geqslant 1$）的升序序列 $S$，处在第 $\lceil L/2 \rceil$ 个位置的数称为 $S$ 的中位数。例如，若序列 $S_1=(11,13,15,17,19)$，则 $S_1$ 的中位数是 15。两个序列的中位数是含它们所有元素的升序序列的中位数。例如，若 $S_2=(2,4,6,8,20)$，则 $S_1$ 和 $S_2$ 的中位数是 11。现有两个等长升序序列 $A$ 和 $B$，设计一个在时间和空间两方面都尽可能高效的算法，找出两个序列 $A$ 和 $B$ 的中位数。

\ifshowsolutions\par\noindent\textbf{解析：}
采用二分法，每次从两个序列中各删除大约一半的元素，直到两个序列均只剩一个元素时取较小者。

\begin{lstlisting}[language=C]
int MidSearch(int A[], int B[], int n) {
    int s1 = 0, d1 = n - 1;  // A 的起止下标
    int s2 = 0, d2 = n - 1;  // B 的起止下标
    int m1, m2;
    while (s1 != d1 || s2 != d2) {
        m1 = (s1 + d1) / 2;
        m2 = (s2 + d2) / 2;
        if (A[m1] == B[m2])
            return A[m1];
        if (A[m1] < B[m2]) {
            // 舍弃 A 的前半段和 B 的后半段
            if ((s1 + d1) % 2 == 0) { s1 = m1; d2 = m2; }
            else                     { s1 = m1 + 1; d2 = m2; }
        } else {
            if ((s2 + d2) % 2 == 0) { d1 = m1; s2 = m2; }
            else                     { d1 = m1; s2 = m2 + 1; }
        }
    }
    return A[s1] < B[s2] ? A[s1] : B[s2];
}
\end{lstlisting}

\textbf{时间复杂度：} $O(\log n)$，每次循环大约减半。\\
\textbf{空间复杂度：} $O(1)$。

\textbf{思想：} 设 $a = A_{\text{med}}$，$b = B_{\text{med}}$。若 $a=b$ 则 $a$ 即为所求；若 $a<b$，则 $A$ 的前半段和 $B$ 的后半段不可能包含中位数，舍弃之；若 $a>b$ 同理。每次舍弃等长元素，保证等长性不变。

\textbf{注意：} 当区间长度为偶数时中位数取偏左位置，元素个数的奇偶性会影响舍弃边界，代码中用 \texttt{(s+d)\%2==0} 判断。
\fi

\item \textbf{（真题风格·综合题）} 已知一个整数序列 $A=(a_0,a_1,\ldots,a_{n-1})$，其中 $0 \leqslant a_i < n$（$0 \leqslant i < n$）。若存在 $a_{p_1}=a_{p_2}=\cdots=a_{p_m}=x$ 且 $m > n/2$（$0 \leqslant p_k < n$，$1 \leqslant k \leqslant m$），则称 $x$ 为 $A$ 的\textbf{主元素}。例如 $A=(0,5,5,3,5,7,5,5)$，则 5 为主元素；又如 $A=(0,5,5,3,5,1,5,7)$，则 $A$ 中没有主元素。假设 $A$ 中的 $n$ 个元素保存在一个一维数组中，设计一个尽可能高效的算法，找出 $A$ 的主元素。若存在主元素，则输出该元素；否则输出 $-1$。

\ifshowsolutions\par\noindent\textbf{解析：}
使用摩尔投票算法（Moore's Voting Algorithm）：候选元素与计数器，一趟扫描确定候选，二趟扫描验证。

\begin{lstlisting}[language=C]
int Majority(int A[], int n) {
    // 第一阶段：选出候选主元素
    int cand = A[0], count = 1;
    for (int i = 1; i < n; i++) {
        if (A[i] == cand) count++;
        else {
            if (count > 0) count--;
            else { cand = A[i]; count = 1; }
        }
    }
    // 第二阶段：验证候选是否确实过半
    if (count == 0) return -1;
    count = 0;
    for (int i = 0; i < n; i++)
        if (A[i] == cand) count++;
    return (count > n / 2) ? cand : -1;
}
\end{lstlisting}

\textbf{时间复杂度：} $O(n)$，两趟线性扫描。\\
\textbf{空间复杂度：} $O(1)$。

\textbf{思想：} 两两抵消——若当前元素与候选相同则计数+1，不同则计数-1。计数归零时更换候选。若存在主元素，最终候选必为主元素（因为主元素出现次数超过半数，无法被完全抵消）。但最终候选不一定过半（如全部不同的序列），故需二趟验证。
\fi

\item \textbf{（真题风格·综合题）} 用单链表保存 $m$ 个整数，结点的结构为 \texttt{[data][next]}，且 $\lvert \texttt{data} \rvert \leqslant n$（$n$ 为正整数）。设计一个时间复杂度尽可能高效的算法，对于链表中 \texttt{data} 的绝对值相等的结点，仅保留第一次出现的结点而删除其余绝对值相等的结点。

\ifshowsolutions\par\noindent\textbf{解析：}
使用辅助数组标记已出现的绝对值，一趟扫描完成。

\begin{lstlisting}[language=C]
void DeleteAbsEqual(LinkList L, int n) {
    int *flag = (int*)calloc(n + 1, sizeof(int)); // flag[x]=1 表示 |x| 已出现
    LNode *p = L, *q;
    while (p->next != NULL) {
        int val = p->next->data;
        int key = (val >= 0) ? val : -val;        // |data|
        if (flag[key] == 1) {                     // 已出现过，删除
            q = p->next;
            p->next = q->next;
            free(q);
        } else {
            flag[key] = 1;
            p = p->next;
        }
    }
    free(flag);
}
\end{lstlisting}

\textbf{时间复杂度：} $O(m)$，遍历链表一次。\\
\textbf{空间复杂度：} $O(n)$，辅助数组大小取决于 \texttt{data} 的取值范围 $n$，而非链表长度。

\textbf{关键：} （1） 单链表删除结点需要知道前驱，用 \texttt{p->next} 指向当前考察结点，避免额外维护前驱指针；（2） 为什么是 $O(n)$ 空间？题目明确 $\lvert \texttt{data} \rvert \leqslant n$，辅助数组大小由 $n$ 决定。408 真题常故意用此考查「时间复杂度尽可能高效」与「空间换时间」的权衡。
\fi

\item \textbf{（真题风格·综合题）} 给定一个含 $n$（$n \geqslant 1$）个整数的数组，设计一个在时间上尽可能高效的算法，找出数组中未出现的最小正整数。例如，数组 $\{-5,3,2,3\}$ 未出现的最小正整数是 1；数组 $\{1,2,3\}$ 未出现的最小正整数是 4。

\ifshowsolutions\par\noindent\textbf{解析：}
核心思想：用原数组本身做哈希表——将每个在 $[1,n]$ 范围内的数放到下标对应的位置。

\begin{lstlisting}[language=C]
int FindMissMin(int A[], int n) {
    // 第一步：将所有在 [1,n] 范围内的数归位
    for (int i = 0; i < n; i++) {
        while (A[i] >= 1 && A[i] <= n && A[A[i] - 1] != A[i]) {
            int temp = A[i];
            A[i] = A[temp - 1];
            A[temp - 1] = temp;
        }
    }
    // 第二步：扫描，第一个 A[i] != i+1 的位置即答案
    for (int i = 0; i < n; i++)
        if (A[i] != i + 1) return i + 1;
    return n + 1;
}
\end{lstlisting}

\textbf{时间复杂度：} $O(n)$。每个元素最多交换一次，\texttt{while} 循环总次数为 $O(n)$。\\
\textbf{空间复杂度：} $O(1)$，原地操作。

\textbf{思想：} 答案只可能在 $[1,n+1]$ 范围内（最坏情况 $1,2,\ldots,n$ 全部出现则答案为 $n+1$）。将每个在 $[1,n]$ 内的数 $x$ 交换到下标 $x-1$ 处，然后检查哪个位置不匹配。
\fi

\item \textbf{（真题风格·综合题）} 已知一个带头结点的单链表，结点结构为 \texttt{[data | next]}。设计一个算法，将该链表按结点位置的奇偶性拆分为两个链表：位置为奇数的结点保持原顺序组成链表 $A$，位置为偶数的结点保持原顺序组成链表 $B$（结点编号从 1 开始计数，即首元结点为第 1 个结点）。

\ifshowsolutions\par\noindent\textbf{解析：}
一趟遍历，用尾插法分别构建两个链表。

\begin{lstlisting}[language=C]
void SplitOddEven(LinkList L, LinkList *A, LinkList *B) {
    *A = (LinkList)malloc(sizeof(LNode));
    *B = (LinkList)malloc(sizeof(LNode));
    (*A)->next = NULL;  (*B)->next = NULL;
    LNode *ra = *A, *rb = *B;      // A, B 的尾指针
    LNode *p = L->next;
    int pos = 1;
    while (p != NULL) {
        LNode *q = p->next;        // 暂存后继
        if (pos % 2 == 1) {        // 奇数位 $\to$ A
            ra->next = p;  ra = p;
        } else {                   // 偶数位 $\to$ B
            rb->next = p;  rb = p;
        }
        p = q;
        pos++;
    }
    ra->next = NULL;  rb->next = NULL;  // 尾结点指针置空
}
\end{lstlisting}

\textbf{时间复杂度：} $O(n)$，一趟遍历。\\
\textbf{空间复杂度：} $O(1)$（仅新建两个头结点，不新建数据结点）。

\textbf{注意：} 拆分后原链表 $L$ 的头结点未释放，调用方根据需要处理。
\fi

\end{enumerate}

\vspace{8pt}
\noindent\textbf{拓展阅读：}
\begin{itemize}
    \item 邓俊辉《数据结构（C++语言版）》第 2 章「向量」和第 3 章「列表」——从 ADT 分层视角理解线性表，包括可扩充向量和有序列表。
    \item Sedgewick《算法》第 1.3 节「背包、队列和栈」——基于链表的实现技巧和迭代器设计。
\end{itemize}
